\section{The common space of spectral relations}
\label{sec:global-defect}

The active space describes the possible responses at one block. We now
identify the relations that remain when many independent blocks are used
together. The essential point is that the derivatives at different
triangular inputs belong to one common space. A bound on their ranks
separately would not give the relations needed for the free-series
construction.

We retain the preceding notation. The allowed sizes are
\(\mathcal S=\{pN:N\geq1\}\); at size \(n\in\mathcal S\), there are
\(e_n=dn/t\) distinct generic roots, each of multiplicity \(t\).
The active space has dimension \(a_n\), and \(\rho_n\) is the generic
rank of the eigenvalue differential, equivalently of the trace
differential. In particular \(\rho_n\leq a_n\).
We use the replication identity \(a_{pN}=Na_p\) and the following
consequence of the preceding attainability argument. There are a size
\(r\in\mathcal S\) and an integer \(K_0\geq1\) such that the loop
responses at a generic size-\(r\) point, together with paired responses
through \(K_0\) independent generic size-\(r\) partners, span its
active space. In each pairing, either arrow may be fixed generically
and the opposite arrow varied. All spaces below are extended to a
common splitting field when several independent points are involved.

\begin{lemma}[Derivatives at a star]
\label{lem:star-derivatives}
Let allowed-size blocks, evaluated by the preceding formal
specialization conventions, have split semisimple outputs with
disjoint spectra and the maximal total number of roots. Add input
arrows only from designated
hub blocks to designated leaf blocks, with no arrows between hubs or
between leaves. The resulting input has the same roots and
multiplicities as its diagonal part. A loop derivative at a block is
its loop response \(D\), with zero entries at all other blocks. If
\(U\) is a forward arrow between a leaf \(P\) and a hub \(Q\), a
derivative in its reverse arrow \(V\) has entries
\[
 \bigl(m_h(P,Q;U,V)\bigr)_h,
 \qquad
 \bigl(m_k(Q,P;V,U)\bigr)_k
\]
at these two blocks and zero entries at the others.
\end{lemma}

\begin{proof}
Order the leaves before the hubs. The input and its evaluated output
are block upper triangular. The diagonal output blocks are the
original ones; since their spectra are disjoint, the output is similar
to their direct sum. In particular it is semisimple with the stated
roots and multiplicities.

Differentiate \(\operatorname{tr}(W^j)\). A closed block path
contributing to a reverse-arrow derivative must use that reverse
arrow once and its matching forward arrow once. There is no other
way to return to the starting block in this star. Its contribution
is therefore exactly the mixed-arrow contribution at the diagonal
input. Loop derivatives use no off-diagonal return. The previously
proved mixed-response formula gives the displayed pairs after
converting trace derivatives to root derivatives. This conversion is
invertible on all distinct roots: its entries are
\(jt\lambda_h^{j-1}\), and the corresponding Vandermonde determinant
is nonzero. The same reasoning applies to specialized blocks whenever
the stated maximal-count and disjointness hypotheses hold.
\end{proof}

\begin{lemma}[A stable upper bound for the differential rank]
\label{lem:tail-defect}
One can choose an arbitrarily large allowed size \(b\) and integers
\(a=a_b\) and \(g\geq0\), such that
\[
 \rho_b=a-g,
 \qquad
 \rho_{bM}\leq Ma-g \quad(M\geq1).
\]
At a generic size-\(b\) point choose one root from each part of the
active partition. Projection onto these \(a\) selected roots
identifies the active space with the full coordinate space. The
quotient of that coordinate space by the loop-response image has
dimension \(g\).
\end{lemma}

\begin{proof}
For \(N\geq1\), put
\[
 g_N=Na_r-\rho_{rN}.
\]
Replication and \(\rho_{rN}\leq a_{rN}\) show that \(g_N\geq0\).
For \(N\geq K_0\), take \(K_0\) of \(N\) independent size-\(r\)
blocks as hubs and the remainder as leaves. By
Lemma~\ref{lem:star-derivatives} and the attainability property, the
projection of the derivative image onto the leaf active spaces is
surjective. Reverse-arrow variables belonging to different leaves
are independent. Thus
\[
 \rho_{rN}\geq(N-K_0)a_r,
 \qquad g_N\leq K_0a_r.
\]
Here the specialized rank is bounded by the generic rank by the
formal minor identities. The finitely many smaller values of \(N\)
do not affect boundedness.

A bounded sequence of nonnegative integers has an eventual minimum
that is attained arbitrarily late. Choose \(N_0\geq K_0\) at such
an occurrence, and write
\[
 g=g_{N_0},\qquad
 g_N\geq g\ (N\geq N_0),\qquad
 b=rN_0,\qquad a=N_0a_r.
\]
Then \(a=a_b\), \(\rho_b=a-g\), and
\[
 \rho_{bM}=Ma-g_{N_0M}\leq Ma-g.
\]
No eventual constancy of \(g_N\) is required. Finally, the active
partition says that projection to one representative of each part
is an isomorphism on the active space. The loop image has rank
\(\rho_b\), so its quotient has dimension \(g\).
\end{proof}

Fix \(b,a,g\) as in Lemma~\ref{lem:tail-defect}. For each generic
size-\(b\) point \(P\), let \(E_P\) be the \(a\)-dimensional
space of selected-root coordinates, and let
\[
 L_P=\operatorname{im}D_P\subseteq E_P,
 \qquad H_P=E_P/L_P.
\]
Thus \(\dim H_P=g\). We regard \(H_P^*\) as the space of
covectors on \(E_P\) annihilating all loop responses.

\begin{lemma}[Surjectivity at both ends of a pair]
\label{lem:pair-quotient-surjectivity}
Let \(P,Q\) be independent generic size-\(b\) points. For a fixed
generic arrow \(U\) between them, varying the reverse arrow \(V\)
gives a joint paired-response map into \(H_P\oplus H_Q\) whose
projection to either summand is surjective. This holds whichever
orientation is fixed.
\end{lemma}

\begin{proof}
Specialize \(P\) and \(Q\) to diagonal sums of \(N_0\)
independent size-\(r\) points. For each little block in \(P\), the
\(N_0\geq K_0\) blocks in \(Q\), together with their components
of the fixed generic arrow, provide enough partners to span its
active space after loop columns are included. The reverse-arrow
entries touching different little blocks of \(P\) vary independently.
Compression at a root of one little block sees only the reverse
entries at that block. Consequently loops and paired columns span
all \(N_0a_r=a\) active coordinates at the \(P\)-end.

The resulting nonzero rank minor persists at generic \(P,Q,U\).
The labeled root data and the pair inverses are regular at the
disjoint diagonal specialization, so this is precisely the
determinant specialization justified earlier. Passing to the loop
quotient gives surjectivity onto \(H_P\). The same proof works at
the \(Q\)-end and with the other orientation fixed. The finitely
many nonzero minor conditions hold simultaneously on a nonempty
generic set.
\end{proof}

The next lemma supplies the common subspace promised at the start
of this section. Its formal-coordinate argument is the reason we
obtain compatible relations between pairs of blocks.

\begin{lemma}[A common tangent space]
\label{lem:common-tangent}
Fix \(M\geq1\) independent generic size-\(b\) points and their
selected roots. At inputs with this diagonal and arbitrary strictly
upper block arrows, the collective span of all selected-root
derivative columns has dimension at most \(Ma-g\). The span of
all loop columns and all two-block paired columns has dimension
exactly \(Ma-g\).
\end{lemma}

\begin{proof}
Put \(R=Ma-g\). For \(M=1\) the assertion is the definition of
the loop rank, so suppose \(M\geq2\). Use one block as a hub,
ordered last. The loops span \(\bigoplus_i L_{P_i}\), of dimension
\(M(a-g)\). Modulo that space, reverse derivatives project
surjectively onto each of the \(M-1\) leaf quotients by
Lemma~\ref{lem:pair-quotient-surjectivity}. These projections vary
independently. The derivative rank at this star is therefore at
least
\[
 M(a-g)+(M-1)g=R.
\]
The generic upper bound in Lemma~\ref{lem:tail-defect} is \(R\).
Thus at a generic strictly upper-arrow array \(\mathcal U\) the
selected-root derivative rank equals \(R\).

Work over the field containing \(\mathcal U\) and the diagonal
splitting data. Let \(F(H)\) be the formal Taylor map of the
selected root shifts under unrestricted input-entry perturbations
\(H\). The base has the maximal root count, so the scalar branch
construction applies. Every \((R+1)\)-minor of \(dF\) is zero,
and some \(R\)-minor is a unit at \(H=0\). Choose the corresponding
\(R\) output coordinates and use them, together with the remaining
input coordinates, as new formal source coordinates \((y,z)\).
The formal inverse theorem gives
\[
 F(y,z)=(y,G(y,z)).
\]
The vanishing \((R+1)\)-minors now imply
\(\partial G/\partial z=0\). In characteristic zero, coefficientwise
differentiation gives \(G(y,z)=G(y)\). At any formal source point
where \(F=0\), all derivative columns therefore belong to the fixed
graph
\[
 \{(v,dG(0)v):v\in E^R\}.
\]

Shifting only the strictly upper arrows leaves every root unchanged.
Apply the preceding conclusion to these shifts: their Jacobian
columns all belong to this one graph, not merely to spaces of the
same dimension.

It remains to pass from formal shifts at \(\mathcal U\) to all
upper-arrow arrays. Let \(J(U)\) be the selected Jacobian at such
an array. It is polynomial in \(U\). Indeed a trace derivative
has just one differentiated input step, and every other
off-diagonal step is strictly upper; a contributing block path
therefore uses a bounded number of upper arrows. The remaining
diagonal sums converge by smallness. The conversion from trace
derivatives to root derivatives has coefficients depending only
on the fixed diagonal roots, so it preserves polynomial dependence.

Every covector annihilating the preceding graph annihilates the
polynomial matrix \(J(\mathcal U+S)\) as a formal identity in
the upper-arrow shift \(S\). Substitution \(S=V-\mathcal U\)
is valid in this polynomial identity. Hence it annihilates every
coefficient of \(J(V)\). The span of those coefficients, formed
over the original diagonal field, has dimension at most \(R\),
since scalar extension preserves rank. It contains every specialized
column and proves the collective-span assertion.

Loop columns occur at the diagonal input. A paired column supported
on two blocks occurs at the base with just the corresponding single
upper arrow, by Lemma~\ref{lem:star-derivatives}. Allowing both
arrows to vary makes the collection independent of which orientation
is called forward. All these columns therefore lie in the common
space. The one-hub star already gives \(R\) independent columns
among them, proving equality.
\end{proof}

\begin{proposition}[Compatible dual relations]
\label{prop:compatible-relations}
For any finite nonempty collection of independent generic size-\(b\)
points \(P_i\), there is a \(g\)-dimensional space
\[
 \mathcal G_{\{P_i\}}\subseteq\bigoplus_i H_{P_i}^*
\]
with the following properties.
\begin{enumerate}
\item Projection to each \(H_{P_i}^*\) is an isomorphism.
\item It annihilates all joint paired responses, with both arrows
arbitrary.
\item For a fixed generic arrow between \(P_i,P_j\), the
annihilator in \(H_{P_i}^*\oplus H_{P_j}^*\) of the variable
reverse-arrow image is exactly the restriction of
\(\mathcal G_{\{P_i\}}\).
\item On adjoining independent blocks, restriction of the larger
space is the original space after scalar extension.
\end{enumerate}
For one block this space is \(H_P^*\). These statements include
\(g=0\).
\end{proposition}

\begin{proof}
Define \(\mathcal G\) as the annihilator of all loop and joint
paired columns in the selected coordinate space. By
Lemma~\ref{lem:common-tangent}, it has dimension \(g\); annihilation
of loops identifies it with a subspace of \(\bigoplus_i H_{P_i}^*\).
If an element has zero component at one vertex \(i\), its pairing
with columns on \((i,j)\) and the surjectivity onto \(H_{P_j}\)
force its component at every other vertex \(j\) to vanish.
Projection onto any one component is therefore injective, and hence
an isomorphism by dimension.

For a fixed generic arrow on \((i,j)\), the reverse-image in
\(H_{P_i}\oplus H_{P_j}\) has dimension at least \(g\), since
its projection onto either summand is surjective. It is killed by
the \(g\)-dimensional restriction of \(\mathcal G\). The latter's
annihilator has dimension \(2g-g=g\), so these two spaces coincide.
This proves the exact pairwise assertion.

Finally, restriction from a larger collection annihilates every old
loop and pair column. It is injective, because projection to any
retained vertex is injective, and both spaces have dimension \(g\).
It is therefore onto. For a single block the same statement follows
from its quotient dimension. All dimension arguments remain valid
when \(g=0\).
\end{proof}

We henceforth write \(\mathcal G\) for these compatible spaces,
specifying the blocks when needed. We have identified all linear
relations shared by the loop and paired corrections. The next
section proves that the leading free-word errors satisfy the
corresponding cyclic relations.
